\documentclass[11pt]{amsart}

\usepackage{amsmath,amssymb,amsthm}
\usepackage[margin=2.6cm]{geometry}
\usepackage{booktabs}
\usepackage{longtable}
\usepackage{seqsplit}
\usepackage[hidelinks]{hyperref}

\newtheorem{theorem}{Theorem}[section]
\newtheorem{lemma}[theorem]{Lemma}
\newtheorem{proposition}[theorem]{Proposition}
\newtheorem{corollary}[theorem]{Corollary}
\theoremstyle{definition}
\newtheorem{definition}[theorem]{Definition}
\theoremstyle{remark}
\newtheorem{remark}[theorem]{Remark}

\newcommand{\lean}[1]{\texttt{\seqsplit{#1}}}
\newcommand{\Mov}{\mathrel{\sim}}
\DeclareMathOperator{\pos}{pos}
\DeclareMathOperator{\out}{out}
\DeclareMathOperator{\len}{len}

\title[J09: word moves as grid moves]{J09: positive word conjugacy as grid moves, and the
component count of word grids}
\date{2 November 2026}

\begin{document}

\begin{abstract}
We prove two statements about the selected grid diagrams of positive braid words (job J09 of
the dependency audit, records M413 and M414). First, if two positive words are related by
\lean{PositiveWordConj}, that is, by rotations, far commutations and positive braid relations
inside words, then their selected grids are related by grid moves. Second, if every generator
occurs in a positive word on $m\ge2$ strands, the number of components of its selected grid equals
the number of cycles, fixed points included, of the permutation of the word. The braid relation is
handled by a finite local model of a window of nine columns, together with a soundness lemma that
transports moves of the local model to grid moves in an arbitrary word.
\end{abstract}

\maketitle

\section{Statements and conventions}\label{sec:statements}

We use the conventions for grid diagrams, cyclic intervals, non-interleaving segments, the moves
of \lean{GridDiagram.IsMove}, and the relation $\Mov$ (\lean{MovesTo}) stated in Section~1 of
\lean{docs/route-n/J11-torus-word-grid.tex}, together with its labeled grid diagrams (Definition~3.1
there) and Lemmas~3.2--3.4 there: an exchange of rows does not change the segments of columns
(Lemma \emph{static}); a column whose two rows are consecutive commutes with every column without
marking in these rows (Lemma \emph{short segments}); and a row with both markings in an interval
$W$ of column positions is non-interleaving with every row with both markings outside $W$ (Lemma
\emph{separated rows}). The destabilization patterns are those described in Section~1.1 there.

\subsection{Word grids}
Let $w$ be a positive word of length $L$ on $m$ strands, with letters
$w_h\in\{0,\dots,m-2\}$ standing for $\sigma_{w_h}$. As in the J10 chapter of
\lean{docs/route-n/written-proofs-J05-J06-J07-J21-J10-J15-J14.tex}, the corners of $w$ are the pairs
$(h,c)$, $0\le h<L$, $c\in\{0,1,2\}$, in the column $3h+c$. The letter $\sigma_s$ in position $h$
involves the ranks $s$ (lower) and $s+1$ (upper). Its incoming corners are $(h,0)$ for the lower and
$(h,1)$ for the upper rank, and its outgoing corners are $(h,1)$ for the lower and $(h,2)$ for the
upper rank. The arc leaving $(h,0)$ has height $s+\frac12$ and ends at $(h,2)$. The arc leaving an
outgoing corner of rank $r$ has height $r$ and ends at the incoming corner of the next occurrence of
$r$ in cyclic order. This defines the permutation $M$ of corners
(\lean{positiveRectilinearWordMatching w}), sending the start of an arc to its end. An
\emph{admissible grid} of $w$ is a grid diagram $G$ of size $3L$ such that, writing
$R(p)=O_G(3h+c)$ for the row of the arc leaving $p=(h,c)$, the $X$-marking of each column $p$ lies in
the row $R(M^{-1}p)$, and $R$ strictly preserves the order of heights
(\lean{IsAdmissibleWordGrid} in \lean{PositiveWordGridAdmissible.lean}). The selected grid
\lean{wordGrid w} is admissible (\lean{wordGrid\_isAdmissible}), and any two admissible grids of $w$
are related by row exchanges inside groups of equal height (\lean{IsAdmissibleWordGrid.movesTo}).
We view an admissible grid as the labeled diagram whose column labels are the corners, in
lexicographic order, and whose row labels are the arcs, identified with their starting corners.

\subsection{The statements}

\begin{theorem}[M413]\label{thm:conj}
Let $u,v$ be positive words on $m$ strands with \lean{PositiveWordConj u v}. Then
$\lean{wordGrid u}\Mov\lean{wordGrid v}$.
\end{theorem}

The audit's acceptance form adds the hypotheses $2\le m$ and that every generator occurs in $u$.
They are not needed.

\begin{theorem}[M414]\label{thm:count}
Let $m\ge2$ and let $w$ be a positive word on $m$ strands in which every generator occurs. Then the
number of components of $\lean{wordGrid w}$ (\lean{componentCount}, the number of cycles of
$c\mapsto X^{-1}(O(c))$) equals \lean{w.toSigned.cycleCount}, the number of orbits of the
permutation of $w$ on $\{0,\dots,m-1\}$, fixed points included.
\end{theorem}

\section{The component count}\label{sec:count}

\subsection{Reduction to the corner matching}
Let $G$ be an admissible grid of $w$ and write $\phi(h,c)=3h+c$. The admissibility equation gives
$X_G(\phi(Mp))=O_G(\phi(p))$ for every corner $p$ (\lean{X\_finProdFinEquiv\_matching}). The
component permutation of $G$ sends $\phi(p)$ to the column of the $X$-marking in the row $O_G(\phi(p))$,
which is $\phi(Mp)$. So the component permutation is $\phi M\phi^{-1}$, and $G$ has as many components
as $M$ has cycles. Every cycle of $M$ is nontrivial, since $M$ has no fixed point
(\lean{positiveRectilinearWordMatching\_ne\_self}). Hence the number of components is the number of
orbits of $M$ (\lean{componentCount\_eq\_orbitCount}).

\subsection{Strand labels}
For $0\le t\le L$ let $P_t$ be the permutation of the prefix of length $t$
(\lean{positiveWordPrefix w t}), so that $P_0=1$, $P_{t+1}=P_t\circ\tau_{w_t}$ with $\tau_s$ the
transposition of $s,s+1$, and $P_L=\pi$ is the permutation of $w$. The value $P_t(r)$ is the initial
position of the strand that lies at position $r$ after $t$ letters. If no letter in positions
$t,\dots,t'-1$ involves the rank $r$, then $P_t(r)=P_{t'}(r)$
(\lean{positiveRectilinearWordPrefix\_constant}).

Define the \emph{label} of a corner by
\[
F(h,0)=F(h,2)=P_h(w_h),\qquad F(h,1)=P_h(w_h+1).
\]
Since $P_{h+1}(w_h)=P_h(w_h+1)$ and $P_{h+1}(w_h+1)=P_h(w_h)$, the label of an outgoing corner of
rank $r$ in position $h$ is $P_{h+1}(r)$, and the label of an incoming corner of rank $r$ in position
$h$ is $P_h(r)$.

\begin{lemma}\label{lem:labels}
Let $p$ be a corner. Exactly one of the following holds.
\begin{enumerate}
\item $\phi(Mp)>\phi(p)$ and $F(Mp)=F(p)$.
\item $p$ is the outgoing corner of rank $r$ of the last occurrence of $r$; then $Mp$ is the incoming
corner of rank $r$ of the first occurrence of $r$, $F(Mp)=r$ and $F(p)=\pi(r)$.
\end{enumerate}
We call the corners of type (2) \emph{wrap corners}.
\end{lemma}

\begin{proof}
If $p=(h,0)$, then $Mp=(h,2)$ and $F(Mp)=F(p)$ by definition. Let $p$ be the outgoing corner of rank
$r$ of an occurrence in position $h$, and let $h'$ be the position of the next occurrence of $r$ in
cyclic order (\lean{positiveRectilinearWordMatching\_occurrence}). If $h'>h$, the corner $Mp$ is in
position $h'$, so $\phi(Mp)\ge 3h'>3h+2\ge\phi(p)$, and no letter in positions $h+1,\dots,h'-1$
involves $r$, so $F(Mp)=P_{h'}(r)=P_{h+1}(r)=F(p)$. Otherwise $h$ is the last occurrence and $h'$ the
first. No letter after $h$ and no letter before $h'$ involves $r$, so $F(p)=P_{h+1}(r)=P_L(r)=\pi(r)$
and $F(Mp)=P_{h'}(r)=P_0(r)=r$. Both statements are the content of
\lean{positiveRectilinearWordMatching\_labels}. In this case $\phi(Mp)\le\phi(p)$: if $h'<h$ this is
clear, and if $h'=h$ the incoming corner of an occurrence precedes its outgoing corner.
\end{proof}

\begin{lemma}\label{lem:reach}
For every corner $p$ there is $n\ge0$ such that $M^np$ is a wrap corner with $F(M^np)=F(p)$. A wrap
corner is determined by its label.
\end{lemma}

\begin{proof}
Induction on $3L-\phi(p)$. If $p$ is a wrap corner, take $n=0$. Otherwise Lemma~\ref{lem:labels}(1)
gives $\phi(Mp)>\phi(p)$ and $F(Mp)=F(p)$, and the induction hypothesis applies to $Mp$. A wrap corner
of rank $r$ has label $\pi(r)$, so its label determines $r$, and the wrap corner of rank $r$ is the
outgoing corner of rank $r$ of the last occurrence of $r$.
\end{proof}

\begin{proposition}\label{prop:orbits}
Assume that every rank occurs in $w$. Then $F$ induces a bijection from the orbits of $M$ onto the
orbits of $\pi$. Precisely, $F$ is surjective, and two corners $p,p'$ lie in the same orbit of $M$ if
and only if $F(p)$ and $F(p')$ lie in the same orbit of $\pi$.
\end{proposition}

\begin{proof}
\emph{Surjectivity.} Let $\ell$ be a position. The rank $\ell$ occurs; let $h$ be its first
occurrence. The incoming corner of rank $\ell$ in position $h$ has label $P_h(\ell)=P_0(\ell)=\ell$.

\emph{Same $M$-orbit implies same $\pi$-orbit.} By Lemma~\ref{lem:labels}, $F(Mp)$ is $F(p)$ or
$\pi^{-1}(F(p))$. By induction, $F(M^np)$ lies in the $\pi$-orbit of $F(p)$ for all $n\ge0$, and every
corner of the $M$-orbit of $p$ has this form because $M$ has finite order.

\emph{Same $\pi$-orbit implies same $M$-orbit.} First let $F(p)=F(p')$. By Lemma~\ref{lem:reach}, $p$
and $p'$ reach wrap corners $z,z'$ with $F(z)=F(p)=F(p')=F(z')$, so $z=z'$, and $p,p'$ lie in one
orbit. In general, let $p$ be a corner and $z$ the wrap corner reached from $p$. By
Lemma~\ref{lem:labels}(2), $F(Mz)=\pi^{-1}(F(z))=\pi^{-1}(F(p))$. By induction on $k$, the $M$-orbit
of $p$ contains a corner of label $\pi^{-k}(F(p))$ for every $k\ge0$. If $F(p')=\pi^j(F(p))$ for some
integer $j$, choose $k\ge0$ with $\pi^{-k}=\pi^j$, which exists because $\pi$ has finite order. The
orbit of $p$ then contains a corner $p''$ with $F(p'')=F(p')$, and $p''$, $p'$ lie in one orbit by the
first case.
\end{proof}

\begin{proof}[Proof of Theorem~\ref{thm:count}]
Since $m\ge2$ and every generator occurs, every rank occurs
(\lean{positiveRectilinearWord\_all\_rows\_occur}). By Proposition~\ref{prop:orbits}, the orbits of $M$
and of $\pi$ are in bijection. The number of orbits of $M$ is the number of components of
$\lean{wordGrid w}$, and the number of orbits of $\pi=\lean{w.toSigned.permutation}$ is
\lean{w.toSigned.cycleCount}.
\end{proof}

\begin{remark}
The hypothesis that every generator occurs cannot be dropped. For $m=3$ and $w=\sigma_0$, the grid
has size $3$ and is a knot, while $\pi$ has the two orbits $\{0,1\}$ and $\{2\}$.
\end{remark}

\section{Rotations and far commutations}\label{sec:easy}

\begin{proposition}\label{prop:rotation}
For every positive word $u$ and letter $i$, $\lean{wordGrid}(u\,i)\Mov\lean{wordGrid}(i\,u)$.
\end{proposition}

\begin{proof}
Let $L$ be the length of $ui$ and $G$ an admissible grid of $ui$. Apply the cyclic permutation of
columns \lean{relabelColumns (finRotate (3L))} three times. The column of the corner $(h,c)$ of $ui$
moves to $3(h+1)+c$ for $h<L-1$ and the columns $3(L-1)+c$ of the last letter move to $c$. Under the
bijection of corners $(h,c)\mapsto(h+1,c)$ for $h<L-1$, $(L-1,c)\mapsto(0,c)$, the corners of $ui$
correspond to the corners of $iu$, letters correspond with equal generators, and the cyclic order of
the occurrences of every rank is rotated. Hence the matching of $ui$ corresponds to the matching of
$iu$, the arcs keep their heights, and the rows are unchanged. The result is an admissible grid of
$iu$, which is related to the selected grid by \lean{IsAdmissibleWordGrid.movesTo}.
\end{proof}

\begin{proposition}\label{prop:far}
If $|i-j|\ge2$, then $\lean{wordGrid}(u\,i\,j\,v)\Mov\lean{wordGrid}(u\,j\,i\,v)$.
\end{proposition}

\begin{proof}
Let $h$ be the length of $u$, and $G$ an admissible grid of $u\,i\,j\,v$. The three columns of the
letter in position $h$ have markings in rows of heights in $\{i,i+\frac12,i+1\}$, and the three
columns of the letter in position $h+1$ in rows of heights in $\{j,j+\frac12,j+1\}$. Since
$|i-j|\ge2$ these sets are separated, and since the rows are ordered by height, the segment of each
column of the first letter is disjoint from the segment of each column of the second letter in the
linear order of rows. Disjoint segments with distinct endpoints are non-interleaving. So the nine
exchanges that move the three columns of position $h+1$ in front of the three columns of position $h$
(keeping their internal order) are grid moves. Under the bijection of corners exchanging the
positions $h$ and $h+1$, the result is an admissible grid of $u\,j\,i\,v$: the two letters involve
disjoint ranks, so the order of the occurrences of each rank is unchanged and the matchings
correspond; the rows are unchanged.
\end{proof}

\section{The braid relation}\label{sec:braid}

Fix $u,v$, a generator $i$ with $i+1\le m-2$, and put $h=\len(u)$,
\[
w=u\,\sigma_i\sigma_{i+1}\sigma_i\,v,\qquad w'=u\,\sigma_{i+1}\sigma_i\sigma_{i+1}\,v.
\]
The \emph{window} is the set $W$ of the nine columns of the letters in positions $h,h+1,h+2$. It is an
interval of column positions in both words. The \emph{block} is the three letters in these positions.

\subsection{Rows of the window}
In an admissible grid of $w$ (or $w'$), a row is \emph{relevant} if it has at least one marking in
$W$, \emph{internal} if both its markings lie in $W$, \emph{connecting} if exactly one does, and
\emph{external} if none does. The relevant rows are the nine arcs leaving the block corners and the
arcs ending at block corners. Only the ranks $i,i+1,i+2$ have integer arcs meeting the block.

For $r\in\{0,1,2\}$, the rank $i+r$ is \emph{degenerate} if it has no occurrence outside the block.
This condition is the same for $w$ and $w'$, since the two words agree outside the block and each
block involves all three ranks. For a nondegenerate rank, the arc ending at the first block occurrence
of $i+r$ comes from outside; we call it $\iota_r$; it is a connecting row with its $O$-marking outside
$W$. The arc leaving the last block occurrence of $i+r$ ends outside; we call it $\theta_r$; it is a
connecting row with its $X$-marking outside $W$. These six rows are the same arcs in the grids of $w$
and of $w'$: $\iota_r$ is the arc leaving the last outside occurrence of $i+r$, and $\theta_r$ is the
arc ending at the first outside occurrence after the block. For a degenerate rank the arc leaving the
last block occurrence ends at the first block occurrence; it is an internal row, which we call
$\zeta_r$ (in $w$) and $\zeta'_r$ (in $w'$). All other relevant rows are internal arcs of the block.
For $\sigma_i\sigma_{i+1}\sigma_i$ they are $\alpha_A,\alpha_B,\alpha_C$ (the middle arcs of the three
letters), $\beta_0$ (rank $i$, from the first to the third letter), $\beta_1$ (rank $i+1$, from the
first to the second letter) and $\beta_2$ (rank $i+1$, from the second to the third letter). For
$\sigma_{i+1}\sigma_i\sigma_{i+1}$ they are $\alpha'_A,\alpha'_B,\alpha'_C$, $\beta'_1$ (rank $i+1$,
first to second letter), $\beta'_2$ (rank $i+2$, first to third letter) and $\beta'_3$ (rank $i+1$,
second to third letter).

\subsection{The local model}
\begin{definition}\label{def:local}
A \emph{window state} consists of a sequence of $k\ge1$ window columns, each with an $O$-label and an
$X$-label, and a sequence of row labels. Each row label is \emph{internal} or \emph{connecting}, and
every connecting label is of type \emph{in} or \emph{out}. Each internal label occurs exactly once as
an $O$-label and once as an $X$-label; each connecting label of type in occurs once as an $X$-label and
never as an $O$-label, and each of type out occurs once as an $O$-label and never as an $X$-label; the
two labels of each column differ. A grid diagram $G$ \emph{realizes} the window state if there is an
interval $W$ of $k$ column positions whose columns carry the given labels, in order, via an injection
of the labels onto rows of $G$ which preserves the order of the sequence of row labels, such that the
internal labels go to rows with both markings in $W$, the labels of type in to rows whose $O$-marking is
outside $W$, the labels of type out to rows whose $X$-marking is outside $W$, and every row of $G$ not
in the image has no marking in $W$.
\end{definition}

For the computation of segments in a window state we place an extra position $k$ to the right of the
window, put the outside marking of every connecting label there, and use cyclic intervals modulo $k+1$.
The \emph{local moves} are:
\begin{enumerate}
\item[(c)] the exchange of two consecutive window columns with non-interleaving segments;
\item[(r)] the exchange of two consecutive row labels, at least one of them internal, with
non-interleaving segments;
\item[(d)] a destabilization whose two columns are consecutive window columns and whose two rows are
consecutive in the sequence of row labels, the deleted row being internal;
\item[(s)] a stabilization at a window column whose new column is a window column and whose new row is
an internal label inserted next to the split row.
\end{enumerate}

\begin{lemma}[soundness]\label{lem:sound}
Let $G$ realize a window state $S$, and let $S'$ be obtained from $S$ by one local move. Then there is a
grid diagram $G'$ realizing $S'$ with $G\Mov G'$, such that $G'$ has the same columns outside the window
with the same markings, the same external rows in the same relative order, and the same relative order
of connecting and external rows.
\end{lemma}

\begin{proof}
(c) The two columns are consecutive columns of $G$. Their markings lie in relevant rows, whose relative
order is the order of the row labels, so the segments in $G$ are non-interleaving if and only if they
are in $S$. The exchange is a grid move.

(r) Let $x$ below $y$ be the two labels, and suppose $x$ is internal. In $G$ the rows strictly between
$x$ and $y$ are external, because the sequence of labels is consecutive. By the lemma on separated rows,
$x$ is non-interleaving with each of them, so $x$ can be moved up until it is immediately below $y$
(Lemma \emph{static}). The segment test for $x,y$ in $G$ agrees with the test in $S$: if $y$ is internal
this is clear; if $y$ is connecting with window marking at position $p$ and outside marking at $q\notin
W$, the cyclic interval $(q,p)^\circ$ meets $W$ in the window positions before $p$ and
$(p,q)^\circ$ meets $W$ in the window positions after $p$, whatever $q$ is, and the cyclic interval
between the two markings of $x$ either is contained in $W$ or contains all positions outside $W$. The
model computes the same sets with $q=k$. So the exchange of $x$ and $y$ is a grid move. If $y$ is
internal and $x$ connecting, move $y$ down instead.

(d) The deleted row is internal. Move it, past external rows only, until it is next to the other row of
the pattern; this is allowed as in (r). The two columns are consecutive in $G$, and the destabilization
pattern only involves these two rows and columns, so it is a destabilization of $G$.

(s) The stabilization of $G$ at the same marking, with the new column next to the split column inside
the window and the new row next to the split row, is an elementary stabilization. The new row is
internal.

In all cases, the external rows and the columns outside $W$ are untouched, except that in (r) and (d)
internal rows move across external rows.
\end{proof}

\subsection{The initial and final windows}
In an admissible grid of $w$ we may choose the order inside groups of equal height. Choose it so that
the relevant rows appear in the order of the \emph{initial window} of $\sigma_i\sigma_{i+1}\sigma_i$
listed in Appendix~\ref{app:tables} for the given set of degenerate ranks. Each listed order is sorted by
height, so this is possible. Then the grid realizes the initial window state. The same applies to $w'$.

Appendix~\ref{app:tables} lists, for each possible set of degenerate ranks, a sequence of local moves
from the initial window of $\sigma_i\sigma_{i+1}\sigma_i$ to a final window, and a sequence from the
initial window of $\sigma_{i+1}\sigma_i\sigma_{i+1}$ to a final window, together with a bijection of the
internal labels of the two final windows which identifies them, columns and row orders included. The sets
$\{i,i+2\}$ cannot occur: if neither $i$ nor $i+2$ occurs outside the block, then neither $\sigma_i$ nor
$\sigma_{i+1}$ occurs outside the block, so $i+1$ does not occur either.

\begin{proposition}\label{prop:braid}
$\lean{wordGrid}(w)\Mov\lean{wordGrid}(w')$.
\end{proposition}

\begin{proof}
Let $G$ and $G'$ be admissible grids of $w$ and $w'$ realizing the two initial windows, with the same
order of the rows that are external or connecting; this is possible because these rows are the same
arcs, with the same heights, in both words. Apply Lemma~\ref{lem:sound} along the two sequences of the
appendix. We obtain $G\Mov H$ and $G'\Mov H'$, where $H$ and $H'$ realize the two final windows. After
the identification of the final windows, $H$ and $H'$ have the same window columns and labels, the same
columns outside the window (the columns of $u$ and $v$, untouched), the same external rows in the same
order, and the same order of connecting and external rows. They can only differ in the positions of
internal rows relative to external rows. Moving internal rows across external rows is allowed by the
lemma on separated rows, so $H\Mov H'$. Hence $G\Mov G'$, and the proposition follows from
\lean{IsAdmissibleWordGrid.movesTo}.
\end{proof}

\begin{remark}
The sequences were found by a computer search in the local model and are listed in full. Each move can
be checked by hand from the listed window states. They were in addition checked by machine in exact
grid diagrams, using the TauCeti definitions of the moves, for six hundred random words containing the
block, covering every set of degenerate ranks that can occur. These checks are evidence only; the proof
is the list of moves together with Lemma~\ref{lem:sound}.
\end{remark}

\section{Proof of Theorem~\ref{thm:conj}}
The relation \lean{PositiveWordConj} is generated by rotations, far commutations and braid relations
inside words, and it is closed under symmetry and transitivity. Propositions~\ref{prop:rotation},
\ref{prop:far} and~\ref{prop:braid} give $\Mov$ for the generators, and $\Mov$ is symmetric and
transitive. Induction on the derivation proves the theorem.

\section{Formalization}\label{sec:formalization}
All results of this note are formalized without additional hypotheses. The component count is
\lean{wordGrid\_componentCount} in \lean{PositiveWordGridComponents.lean}. The rotations and the far
commutations are \lean{wordGrid\_movesTo\_rotate} and \lean{wordGrid\_movesTo\_far} in
\lean{PositiveWordGridMovesTransport.lean}. Proposition~\ref{prop:braid}, for all words $u,v$ and all
adjacent letters $i$ and $j=i+1$, is \lean{wordGrid\_movesTo\_braid'} in
\lean{PositiveWordGridBraidRelationFront.lean}, and in the form \lean{WordGridBraidRelationStatement}
it is \lean{wordGrid\_movesTo\_braid} in \lean{PositiveWordGridBraidRelation.lean}.
Together with \lean{wordGrid\_movesTo\_of\_conj\_of\_braid} in \lean{PositiveWordGridMoves.lean} it
gives Theorem~\ref{thm:conj} (\lean{BraidWindow.wordGridConjugacyMovesStatement\_holds}).

The formal proof of Proposition~\ref{prop:braid} follows Section~\ref{sec:braid}, with the following
choices.
\begin{enumerate}
\item The block is moved to the front of the word by rotations
(\lean{wordGrid\_movesTo\_append\_comm}), so the window consists of the first nine columns and the
rest of the word is $v\,u$.
\item A window state is a list of label pairs for the window columns and a list of row labels, with
labels encoded by natural numbers (\lean{PositiveWordGridBraidRelationModel.lean}). The function
\lean{BraidWindow.run} applies a list of local moves and checks their conditions. A realization
(\lean{BraidWindow.Real}) is a labeled grid whose first columns carry the window labels and whose
window rows appear in the order of the window state. Lemma~\ref{lem:sound}, iterated along a list of
moves, is \lean{BraidWindow.run\_sound}.
\item The initial grid \lean{BraidWindow.initGrid} has the corners of the word as column and row
labels, its rows sorted by a key that refines the heights of the arcs, and it is relabeled by the window
labels. It realizes the initial window (\lean{BraidWindow.initGrid\_real}) and is related to the
selected grid (\lean{BraidWindow.initGrid\_movesTo}). The matching at the block corners is computed by
\lean{BraidWindow.BM} (\lean{BraidWindow.BlockHyp.blockM\_some},
\lean{BraidWindow.BlockHyp.blockM\_none}). The two initial grids agree outside the window
(\lean{BraidWindow.agree\_O}, \lean{BraidWindow.agree\_X}, \lean{BraidWindow.agree\_rows}).
\item The data of Appendix~\ref{app:tables} is in \lean{PositiveWordGridBraidRelationData.lean}. For
each set of degenerate ranks, the conditions on the initial windows (\lean{BraidWindow.InitOK},
\lean{BraidWindow.PairOK}), the two sequences of moves and the identification of the final windows
(\lean{BraidWindow.FinalOK}) are proved by kernel evaluation with \lean{decide}. These checks are part
of the formal proof, unlike the random checks mentioned in the remark after
Proposition~\ref{prop:braid}.
\item The final comparison of the two final grids is \lean{BraidWindow.wordGrid\_movesTo\_of\_windows}.
It moves internal rows across external rows by \lean{LabeledGrid.movesTo\_of\_shuffle}.
\end{enumerate}

\appendix
\section{The local sequences}\label{app:tables}
\sloppy

Notation. Window columns are numbered from $0$ in their current order, and each is written as the pair
$(O\text{-label},X\text{-label})$. Rows are listed from bottom to top. The moves are:
$\mathsf c_t$, the exchange of columns $t$ and $t+1$; $\mathsf r(x,y)$, the exchange of the consecutive
rows $x$ (below) and $y$ (above); $\mathsf d^{X}(p,s;x\to y)$, the $X$-destabilization in which the
corners $(p,x),(s,x),(p,y)$ carry $O,X,X$ and $(s,y)$ carries no $O$, deleting column $p$ and row $x$
and moving the $X$-marking of column $s$ to row $y$; $\mathsf d^{O}(p,s;x\to y)$, the same with $O$ and
$X$ exchanged; and $\mathsf s^{O}(s,\mathrm{L/R},\mathrm{b/a};\nu)$, the $O$-stabilization splitting the
$O$-marking of column $s$, lying in row $y$, with the new column immediately left (L) or right (R) of
column $s$ and the new row $\nu$ immediately below (b) or above (a) row $y$, after which column $s$ has
$O$-label $\nu$ and the new column has labels $(y,\nu)$; $\mathsf s^{X}$ is the same with $O$ and $X$
exchanged. The connecting labels are $\iota_r$ (type in) and $\theta_r$ (type out); all other labels are
internal. The labels $\nu_n$ are new internal rows.

\subsection*{Case: no degenerate rank}
\noindent Initial window of $\sigma_i\sigma_{i+1}\sigma_i$. Columns: $(\alpha_A,\iota_0),\allowbreak (\beta_0,\iota_1),\allowbreak (\beta_1,\alpha_A),\allowbreak (\alpha_B,\beta_1),\allowbreak (\beta_2,\iota_2),\allowbreak (\theta_2,\alpha_B),\allowbreak (\alpha_C,\beta_0),\allowbreak (\theta_0,\beta_2),\allowbreak (\theta_1,\alpha_C)$.\\ Rows: $\theta_0 <\allowbreak \iota_0 <\allowbreak \beta_0 <\allowbreak \alpha_A <\allowbreak \alpha_C <\allowbreak \iota_1 <\allowbreak \beta_1 <\allowbreak \beta_2 <\allowbreak \theta_1 <\allowbreak \alpha_B <\allowbreak \theta_2 <\allowbreak \iota_2$.\par
\noindent Moves (17): $\mathsf{r}(\alpha_C,\iota_1)$, $\mathsf{c}_{4}$, $\mathsf{d}^{O}(4,3;\alpha_B\to\theta_2)$, $\mathsf{r}(\iota_0,\beta_0)$, $\mathsf{c}_{0}$, $\mathsf{d}^{X}(1,2;\alpha_A\to\iota_0)$, $\mathsf{r}(\alpha_C,\beta_1)$, $\mathsf{r}(\beta_2,\theta_1)$, $\mathsf{r}(\iota_1,\beta_1)$, $\mathsf{d}^{X}(1,2;\beta_1\to\iota_0)$, $\mathsf{r}(\beta_2,\theta_2)$, $\mathsf{c}_{3}$, $\mathsf{d}^{O}(5,4;\alpha_C\to\theta_1)$, $\mathsf{d}^{X}(2,3;\beta_2\to\iota_2)$, $\mathsf{c}_{1}$, $\mathsf{s}^{O}(3,R,a;\nu_{516})$, $\mathsf{r}(\nu_{516},\theta_2)$.\par
\noindent Final window. Columns: $(\beta_0,\iota_1),\allowbreak (\theta_0,\iota_2),\allowbreak (\theta_2,\iota_0),\allowbreak (\nu_{516},\beta_0),\allowbreak (\theta_1,\nu_{516})$.\\ Rows: $\theta_0 <\allowbreak \beta_0 <\allowbreak \iota_0 <\allowbreak \iota_1 <\allowbreak \theta_1 <\allowbreak \theta_2 <\allowbreak \nu_{516} <\allowbreak \iota_2$.\par\medskip
\noindent Initial window of $\sigma_{i+1}\sigma_i\sigma_{i+1}$. Columns: $(\alpha'_A,\iota_1),\allowbreak (\beta'_1,\iota_2),\allowbreak (\beta'_2,\alpha'_A),\allowbreak (\alpha'_B,\iota_0),\allowbreak (\theta_0,\beta'_1),\allowbreak (\beta'_3,\alpha'_B),\allowbreak (\alpha'_C,\beta'_3),\allowbreak (\theta_1,\beta'_2),\allowbreak (\theta_2,\alpha'_C)$.\\ Rows: $\theta_0 <\allowbreak \iota_0 <\allowbreak \alpha'_B <\allowbreak \iota_1 <\allowbreak \beta'_1 <\allowbreak \beta'_3 <\allowbreak \theta_1 <\allowbreak \alpha'_A <\allowbreak \alpha'_C <\allowbreak \theta_2 <\allowbreak \beta'_2 <\allowbreak \iota_2$.\par
\noindent Moves (19): $\mathsf{r}(\iota_1,\beta'_1)$, $\mathsf{c}_{3}$, $\mathsf{d}^{X}(4,5;\alpha'_B\to\iota_0)$, $\mathsf{r}(\iota_1,\beta'_3)$, $\mathsf{r}(\beta'_1,\beta'_3)$, $\mathsf{d}^{X}(4,5;\beta'_3\to\iota_0)$, $\mathsf{r}(\iota_0,\beta'_1)$, $\mathsf{c}_{1}$, $\mathsf{d}^{O}(3,2;\beta'_1\to\theta_0)$, $\mathsf{r}(\alpha'_A,\alpha'_C)$, $\mathsf{r}(\alpha'_A,\theta_2)$, $\mathsf{d}^{O}(1,0;\alpha'_A\to\beta'_2)$, $\mathsf{c}_{0}$, $\mathsf{c}_{3}$, $\mathsf{d}^{O}(3,2;\alpha'_C\to\theta_2)$, $\mathsf{s}^{X}(1,L,b;\nu_{716})$, $\mathsf{r}(\iota_0,\nu_{716})$, $\mathsf{c}_{0}$, $\mathsf{c}_{2}$.\par
\noindent Final window. Columns: $(\nu_{716},\iota_1),\allowbreak (\theta_0,\iota_2),\allowbreak (\theta_2,\iota_0),\allowbreak (\beta'_2,\nu_{716}),\allowbreak (\theta_1,\beta'_2)$.\\ Rows: $\theta_0 <\allowbreak \nu_{716} <\allowbreak \iota_0 <\allowbreak \iota_1 <\allowbreak \theta_1 <\allowbreak \theta_2 <\allowbreak \beta'_2 <\allowbreak \iota_2$.\par\medskip
\noindent Identification of the final windows: $\beta_0\leftrightarrow \nu_{716},\allowbreak \\nu_{516}\leftrightarrow \beta'_2$.\par\bigskip
\subsection*{Case: degenerate ranks \texorpdfstring{$\{i\}$}{i}}
\noindent Initial window of $\sigma_i\sigma_{i+1}\sigma_i$. Columns: $(\alpha_A,\zeta_0),\allowbreak (\beta_0,\iota_1),\allowbreak (\beta_1,\alpha_A),\allowbreak (\alpha_B,\beta_1),\allowbreak (\beta_2,\iota_2),\allowbreak (\theta_2,\alpha_B),\allowbreak (\alpha_C,\beta_0),\allowbreak (\zeta_0,\beta_2),\allowbreak (\theta_1,\alpha_C)$.\\ Rows: $\zeta_0 <\allowbreak \beta_0 <\allowbreak \alpha_A <\allowbreak \alpha_C <\allowbreak \iota_1 <\allowbreak \beta_1 <\allowbreak \beta_2 <\allowbreak \theta_1 <\allowbreak \alpha_B <\allowbreak \theta_2 <\allowbreak \iota_2$.\par
\noindent Moves (23): $\mathsf{r}(\alpha_C,\iota_1)$, $\mathsf{c}_{4}$, $\mathsf{d}^{O}(4,3;\alpha_B\to\theta_2)$, $\mathsf{r}(\zeta_0,\beta_0)$, $\mathsf{c}_{0}$, $\mathsf{d}^{X}(1,2;\alpha_A\to\zeta_0)$, $\mathsf{r}(\iota_1,\alpha_C)$, $\mathsf{r}(\beta_1,\beta_2)$, $\mathsf{r}(\beta_1,\theta_1)$, $\mathsf{d}^{O}(2,1;\beta_1\to\theta_2)$, $\mathsf{r}(\alpha_C,\iota_1)$, $\mathsf{r}(\zeta_0,\iota_1)$, $\mathsf{c}_{0}$, $\mathsf{c}_{2}$, $\mathsf{d}^{X}(1,2;\beta_0\to\iota_1)$, $\mathsf{r}(\beta_2,\theta_1)$, $\mathsf{r}(\alpha_C,\theta_1)$, $\mathsf{r}(\zeta_0,\theta_1)$, $\mathsf{r}(\alpha_C,\beta_2)$, $\mathsf{d}^{O}(3,2;\beta_2\to\zeta_0)$, $\mathsf{r}(\alpha_C,\theta_2)$, $\mathsf{c}_{0}$, $\mathsf{d}^{O}(1,2;\zeta_0\to\theta_2)$.\par
\noindent Final window. Columns: $(\alpha_C,\iota_1),\allowbreak (\theta_2,\iota_2),\allowbreak (\theta_1,\alpha_C)$.\\ Rows: $\iota_1 <\allowbreak \theta_1 <\allowbreak \theta_2 <\allowbreak \alpha_C <\allowbreak \iota_2$.\par\medskip
\noindent Initial window of $\sigma_{i+1}\sigma_i\sigma_{i+1}$. Columns: $(\alpha'_A,\iota_1),\allowbreak (\beta'_1,\iota_2),\allowbreak (\beta'_2,\alpha'_A),\allowbreak (\alpha'_B,\zeta'_0),\allowbreak (\zeta'_0,\beta'_1),\allowbreak (\beta'_3,\alpha'_B),\allowbreak (\alpha'_C,\beta'_3),\allowbreak (\theta_1,\beta'_2),\allowbreak (\theta_2,\alpha'_C)$.\\ Rows: $\zeta'_0 <\allowbreak \alpha'_B <\allowbreak \iota_1 <\allowbreak \beta'_1 <\allowbreak \beta'_3 <\allowbreak \theta_1 <\allowbreak \alpha'_A <\allowbreak \alpha'_C <\allowbreak \theta_2 <\allowbreak \beta'_2 <\allowbreak \iota_2$.\par
\noindent Moves (14): $\mathsf{d}^{O}(3,4;\zeta'_0\to\alpha'_B)$, $\mathsf{r}(\beta'_3,\theta_1)$, $\mathsf{r}(\alpha'_B,\iota_1)$, $\mathsf{d}^{X}(3,4;\alpha'_B\to\beta'_1)$, $\mathsf{r}(\theta_1,\beta'_3)$, $\mathsf{d}^{X}(3,4;\beta'_3\to\beta'_1)$, $\mathsf{r}(\beta'_1,\theta_1)$, $\mathsf{c}_{3}$, $\mathsf{d}^{O}(5,4;\alpha'_C\to\theta_2)$, $\mathsf{c}_{1}$, $\mathsf{r}(\alpha'_A,\theta_2)$, $\mathsf{d}^{O}(1,0;\alpha'_A\to\beta'_2)$, $\mathsf{c}_{2}$, $\mathsf{d}^{O}(2,1;\beta'_1\to\theta_2)$.\par
\noindent Final window. Columns: $(\beta'_2,\iota_1),\allowbreak (\theta_2,\iota_2),\allowbreak (\theta_1,\beta'_2)$.\\ Rows: $\iota_1 <\allowbreak \theta_1 <\allowbreak \theta_2 <\allowbreak \beta'_2 <\allowbreak \iota_2$.\par\medskip
\noindent Identification of the final windows: $\alpha_C\leftrightarrow \beta'_2$.\par\bigskip
\subsection*{Case: degenerate ranks \texorpdfstring{$\{i+1\}$}{i+1}}
\noindent Initial window of $\sigma_i\sigma_{i+1}\sigma_i$. Columns: $(\alpha_A,\iota_0),\allowbreak (\beta_0,\zeta_1),\allowbreak (\beta_1,\alpha_A),\allowbreak (\alpha_B,\beta_1),\allowbreak (\beta_2,\iota_2),\allowbreak (\theta_2,\alpha_B),\allowbreak (\alpha_C,\beta_0),\allowbreak (\theta_0,\beta_2),\allowbreak (\zeta_1,\alpha_C)$.\\ Rows: $\theta_0 <\allowbreak \iota_0 <\allowbreak \beta_0 <\allowbreak \alpha_A <\allowbreak \alpha_C <\allowbreak \zeta_1 <\allowbreak \beta_1 <\allowbreak \beta_2 <\allowbreak \alpha_B <\allowbreak \theta_2 <\allowbreak \iota_2$.\par
\noindent Moves (16): $\mathsf{r}(\iota_0,\beta_0)$, $\mathsf{c}_{4}$, $\mathsf{d}^{O}(4,3;\alpha_B\to\theta_2)$, $\mathsf{c}_{5}$, $\mathsf{d}^{O}(7,6;\alpha_C\to\zeta_1)$, $\mathsf{c}_{4}$, $\mathsf{r}(\zeta_1,\beta_1)$, $\mathsf{d}^{X}(2,3;\beta_1\to\alpha_A)$, $\mathsf{r}(\beta_2,\theta_2)$, $\mathsf{d}^{X}(4,3;\beta_2\to\iota_2)$, $\mathsf{c}_{0}$, $\mathsf{d}^{X}(1,2;\alpha_A\to\iota_0)$, $\mathsf{c}_{1}$, $\mathsf{s}^{X}(0,L,a;\nu_{514})$, $\mathsf{r}(\nu_{514},\theta_2)$, $\mathsf{c}_{1}$.\par
\noindent Final window. Columns: $(\nu_{514},\zeta_1),\allowbreak (\theta_0,\iota_2),\allowbreak (\beta_0,\nu_{514}),\allowbreak (\theta_2,\iota_0),\allowbreak (\zeta_1,\beta_0)$.\\ Rows: $\theta_0 <\allowbreak \beta_0 <\allowbreak \iota_0 <\allowbreak \zeta_1 <\allowbreak \theta_2 <\allowbreak \nu_{514} <\allowbreak \iota_2$.\par\medskip
\noindent Initial window of $\sigma_{i+1}\sigma_i\sigma_{i+1}$. Columns: $(\alpha'_A,\zeta'_1),\allowbreak (\beta'_1,\iota_2),\allowbreak (\beta'_2,\alpha'_A),\allowbreak (\alpha'_B,\iota_0),\allowbreak (\theta_0,\beta'_1),\allowbreak (\beta'_3,\alpha'_B),\allowbreak (\alpha'_C,\beta'_3),\allowbreak (\zeta'_1,\beta'_2),\allowbreak (\theta_2,\alpha'_C)$.\\ Rows: $\theta_0 <\allowbreak \iota_0 <\allowbreak \alpha'_B <\allowbreak \zeta'_1 <\allowbreak \beta'_1 <\allowbreak \beta'_3 <\allowbreak \alpha'_A <\allowbreak \alpha'_C <\allowbreak \theta_2 <\allowbreak \beta'_2 <\allowbreak \iota_2$.\par
\noindent Moves (17): $\mathsf{c}_{6}$, $\mathsf{d}^{O}(8,7;\alpha'_C\to\theta_2)$, $\mathsf{c}_{3}$, $\mathsf{d}^{X}(4,5;\alpha'_B\to\iota_0)$, $\mathsf{r}(\zeta'_1,\beta'_1)$, $\mathsf{r}(\iota_0,\beta'_1)$, $\mathsf{c}_{2}$, $\mathsf{d}^{O}(2,1;\beta'_1\to\theta_0)$, $\mathsf{c}_{4}$, $\mathsf{r}(\zeta'_1,\beta'_3)$, $\mathsf{d}^{X}(3,4;\beta'_3\to\iota_0)$, $\mathsf{r}(\alpha'_A,\theta_2)$, $\mathsf{c}_{2}$, $\mathsf{d}^{X}(3,4;\beta'_2\to\alpha'_A)$, $\mathsf{s}^{O}(3,R,b;\nu_{715})$, $\mathsf{r}(\iota_0,\nu_{715})$, $\mathsf{c}_{2}$.\par
\noindent Final window. Columns: $(\alpha'_A,\zeta'_1),\allowbreak (\theta_0,\iota_2),\allowbreak (\nu_{715},\alpha'_A),\allowbreak (\theta_2,\iota_0),\allowbreak (\zeta'_1,\nu_{715})$.\\ Rows: $\theta_0 <\allowbreak \nu_{715} <\allowbreak \iota_0 <\allowbreak \zeta'_1 <\allowbreak \theta_2 <\allowbreak \alpha'_A <\allowbreak \iota_2$.\par\medskip
\noindent Identification of the final windows: $\beta_0\leftrightarrow \nu_{715},\allowbreak \\zeta_1\leftrightarrow \zeta'_1,\allowbreak \\nu_{514}\leftrightarrow \alpha'_A$.\par\bigskip
\subsection*{Case: degenerate ranks \texorpdfstring{$\{i+2\}$}{i+2}}
\noindent Initial window of $\sigma_i\sigma_{i+1}\sigma_i$. Columns: $(\alpha_A,\iota_0),\allowbreak (\beta_0,\iota_1),\allowbreak (\beta_1,\alpha_A),\allowbreak (\alpha_B,\beta_1),\allowbreak (\beta_2,\zeta_2),\allowbreak (\zeta_2,\alpha_B),\allowbreak (\alpha_C,\beta_0),\allowbreak (\theta_0,\beta_2),\allowbreak (\theta_1,\alpha_C)$.\\ Rows: $\theta_0 <\allowbreak \iota_0 <\allowbreak \beta_0 <\allowbreak \alpha_A <\allowbreak \alpha_C <\allowbreak \iota_1 <\allowbreak \beta_1 <\allowbreak \beta_2 <\allowbreak \theta_1 <\allowbreak \alpha_B <\allowbreak \zeta_2$.\par
\noindent Moves (15): $\mathsf{d}^{X}(5,4;\zeta_2\to\alpha_B)$, $\mathsf{r}(\iota_0,\beta_0)$, $\mathsf{c}_{0}$, $\mathsf{d}^{X}(1,2;\alpha_A\to\iota_0)$, $\mathsf{r}(\iota_0,\alpha_C)$, $\mathsf{c}_{4}$, $\mathsf{d}^{X}(5,6;\alpha_C\to\beta_0)$, $\mathsf{c}_{3}$, $\mathsf{r}(\iota_1,\beta_1)$, $\mathsf{d}^{X}(1,2;\beta_1\to\iota_0)$, $\mathsf{c}_{2}$, $\mathsf{r}(\theta_1,\alpha_B)$, $\mathsf{d}^{O}(2,1;\alpha_B\to\beta_2)$, $\mathsf{r}(\iota_1,\beta_2)$, $\mathsf{d}^{X}(1,2;\beta_2\to\iota_0)$.\par
\noindent Final window. Columns: $(\beta_0,\iota_1),\allowbreak (\theta_0,\iota_0),\allowbreak (\theta_1,\beta_0)$.\\ Rows: $\theta_0 <\allowbreak \beta_0 <\allowbreak \iota_0 <\allowbreak \iota_1 <\allowbreak \theta_1$.\par\medskip
\noindent Initial window of $\sigma_{i+1}\sigma_i\sigma_{i+1}$. Columns: $(\alpha'_A,\iota_1),\allowbreak (\beta'_1,\zeta'_2),\allowbreak (\beta'_2,\alpha'_A),\allowbreak (\alpha'_B,\iota_0),\allowbreak (\theta_0,\beta'_1),\allowbreak (\beta'_3,\alpha'_B),\allowbreak (\alpha'_C,\beta'_3),\allowbreak (\theta_1,\beta'_2),\allowbreak (\zeta'_2,\alpha'_C)$.\\ Rows: $\theta_0 <\allowbreak \iota_0 <\allowbreak \alpha'_B <\allowbreak \iota_1 <\allowbreak \beta'_1 <\allowbreak \beta'_3 <\allowbreak \theta_1 <\allowbreak \alpha'_A <\allowbreak \alpha'_C <\allowbreak \zeta'_2 <\allowbreak \beta'_2$.\par
\noindent Moves (19): $\mathsf{c}_{3}$, $\mathsf{d}^{X}(4,5;\alpha'_B\to\iota_0)$, $\mathsf{r}(\iota_1,\beta'_1)$, $\mathsf{c}_{6}$, $\mathsf{d}^{O}(6,5;\alpha'_C\to\zeta'_2)$, $\mathsf{r}(\iota_1,\beta'_3)$, $\mathsf{r}(\beta'_1,\beta'_3)$, $\mathsf{d}^{X}(4,5;\beta'_3\to\iota_0)$, $\mathsf{r}(\iota_0,\beta'_1)$, $\mathsf{r}(\theta_1,\alpha'_A)$, $\mathsf{c}_{2}$, $\mathsf{d}^{O}(2,1;\beta'_1\to\theta_0)$, $\mathsf{r}(\iota_1,\alpha'_A)$, $\mathsf{r}(\iota_0,\alpha'_A)$, $\mathsf{r}(\theta_1,\zeta'_2)$, $\mathsf{c}_{2}$, $\mathsf{d}^{O}(4,3;\beta'_2\to\theta_1)$, $\mathsf{r}(\iota_1,\zeta'_2)$, $\mathsf{d}^{X}(2,1;\zeta'_2\to\iota_0)$.\par
\noindent Final window. Columns: $(\alpha'_A,\iota_1),\allowbreak (\theta_0,\iota_0),\allowbreak (\theta_1,\alpha'_A)$.\\ Rows: $\theta_0 <\allowbreak \alpha'_A <\allowbreak \iota_0 <\allowbreak \iota_1 <\allowbreak \theta_1$.\par\medskip
\noindent Identification of the final windows: $\beta_0\leftrightarrow \alpha'_A$.\par\bigskip
\subsection*{Case: degenerate ranks \texorpdfstring{$\{i,i+1\}$}{i,i+1}}
\noindent Initial window of $\sigma_i\sigma_{i+1}\sigma_i$. Columns: $(\alpha_A,\zeta_0),\allowbreak (\beta_0,\zeta_1),\allowbreak (\beta_1,\alpha_A),\allowbreak (\alpha_B,\beta_1),\allowbreak (\beta_2,\iota_2),\allowbreak (\theta_2,\alpha_B),\allowbreak (\alpha_C,\beta_0),\allowbreak (\zeta_0,\beta_2),\allowbreak (\zeta_1,\alpha_C)$.\\ Rows: $\zeta_0 <\allowbreak \beta_0 <\allowbreak \alpha_A <\allowbreak \alpha_C <\allowbreak \zeta_1 <\allowbreak \beta_1 <\allowbreak \beta_2 <\allowbreak \alpha_B <\allowbreak \theta_2 <\allowbreak \iota_2$.\par
\noindent Moves (19): $\mathsf{c}_{4}$, $\mathsf{d}^{O}(4,3;\alpha_B\to\theta_2)$, $\mathsf{r}(\zeta_0,\beta_0)$, $\mathsf{r}(\beta_1,\beta_2)$, $\mathsf{d}^{O}(3,2;\beta_1\to\theta_2)$, $\mathsf{r}(\zeta_0,\alpha_A)$, $\mathsf{c}_{0}$, $\mathsf{d}^{X}(1,2;\alpha_A\to\zeta_0)$, $\mathsf{c}_{2}$, $\mathsf{c}_{3}$, $\mathsf{r}(\beta_2,\theta_2)$, $\mathsf{d}^{X}(4,3;\beta_2\to\iota_2)$, $\mathsf{r}(\alpha_C,\zeta_1)$, $\mathsf{r}(\zeta_0,\zeta_1)$, $\mathsf{r}(\alpha_C,\theta_2)$, $\mathsf{c}_{0}$, $\mathsf{d}^{X}(1,2;\beta_0\to\zeta_1)$, $\mathsf{c}_{0}$, $\mathsf{d}^{O}(1,2;\zeta_0\to\theta_2)$.\par
\noindent Final window. Columns: $(\alpha_C,\zeta_1),\allowbreak (\theta_2,\iota_2),\allowbreak (\zeta_1,\alpha_C)$.\\ Rows: $\zeta_1 <\allowbreak \theta_2 <\allowbreak \alpha_C <\allowbreak \iota_2$.\par\medskip
\noindent Initial window of $\sigma_{i+1}\sigma_i\sigma_{i+1}$. Columns: $(\alpha'_A,\zeta'_1),\allowbreak (\beta'_1,\iota_2),\allowbreak (\beta'_2,\alpha'_A),\allowbreak (\alpha'_B,\zeta'_0),\allowbreak (\zeta'_0,\beta'_1),\allowbreak (\beta'_3,\alpha'_B),\allowbreak (\alpha'_C,\beta'_3),\allowbreak (\zeta'_1,\beta'_2),\allowbreak (\theta_2,\alpha'_C)$.\\ Rows: $\zeta'_0 <\allowbreak \alpha'_B <\allowbreak \zeta'_1 <\allowbreak \beta'_1 <\allowbreak \beta'_3 <\allowbreak \alpha'_A <\allowbreak \alpha'_C <\allowbreak \theta_2 <\allowbreak \beta'_2 <\allowbreak \iota_2$.\par
\noindent Moves (13): $\mathsf{d}^{O}(3,4;\zeta'_0\to\alpha'_B)$, $\mathsf{c}_{5}$, $\mathsf{d}^{O}(7,6;\alpha'_C\to\theta_2)$, $\mathsf{r}(\alpha'_A,\theta_2)$, $\mathsf{r}(\alpha'_B,\zeta'_1)$, $\mathsf{d}^{X}(3,4;\alpha'_B\to\beta'_1)$, $\mathsf{r}(\beta'_1,\beta'_3)$, $\mathsf{c}_{3}$, $\mathsf{d}^{X}(2,3;\beta'_2\to\alpha'_A)$, $\mathsf{d}^{X}(3,4;\beta'_3\to\beta'_1)$, $\mathsf{r}(\beta'_1,\theta_2)$, $\mathsf{c}_{2}$, $\mathsf{d}^{O}(2,1;\beta'_1\to\theta_2)$.\par
\noindent Final window. Columns: $(\alpha'_A,\zeta'_1),\allowbreak (\theta_2,\iota_2),\allowbreak (\zeta'_1,\alpha'_A)$.\\ Rows: $\zeta'_1 <\allowbreak \theta_2 <\allowbreak \alpha'_A <\allowbreak \iota_2$.\par\medskip
\noindent Identification of the final windows: $\zeta_1\leftrightarrow \zeta'_1,\allowbreak \\alpha_C\leftrightarrow \alpha'_A$.\par\bigskip
\subsection*{Case: degenerate ranks \texorpdfstring{$\{i+1,i+2\}$}{i+1,i+2}}
\noindent Initial window of $\sigma_i\sigma_{i+1}\sigma_i$. Columns: $(\alpha_A,\iota_0),\allowbreak (\beta_0,\zeta_1),\allowbreak (\beta_1,\alpha_A),\allowbreak (\alpha_B,\beta_1),\allowbreak (\beta_2,\zeta_2),\allowbreak (\zeta_2,\alpha_B),\allowbreak (\alpha_C,\beta_0),\allowbreak (\theta_0,\beta_2),\allowbreak (\zeta_1,\alpha_C)$.\\ Rows: $\theta_0 <\allowbreak \iota_0 <\allowbreak \beta_0 <\allowbreak \alpha_A <\allowbreak \alpha_C <\allowbreak \zeta_1 <\allowbreak \beta_1 <\allowbreak \beta_2 <\allowbreak \alpha_B <\allowbreak \zeta_2$.\par
\noindent Moves (12): $\mathsf{d}^{X}(5,4;\zeta_2\to\alpha_B)$, $\mathsf{d}^{O}(4,3;\alpha_B\to\beta_2)$, $\mathsf{d}^{O}(3,2;\beta_1\to\beta_2)$, $\mathsf{c}_{4}$, $\mathsf{d}^{O}(4,3;\alpha_C\to\zeta_1)$, $\mathsf{r}(\iota_0,\beta_0)$, $\mathsf{r}(\iota_0,\alpha_A)$, $\mathsf{c}_{0}$, $\mathsf{d}^{X}(1,2;\alpha_A\to\iota_0)$, $\mathsf{c}_{2}$, $\mathsf{r}(\zeta_1,\beta_2)$, $\mathsf{d}^{X}(1,2;\beta_2\to\iota_0)$.\par
\noindent Final window. Columns: $(\beta_0,\zeta_1),\allowbreak (\theta_0,\iota_0),\allowbreak (\zeta_1,\beta_0)$.\\ Rows: $\theta_0 <\allowbreak \beta_0 <\allowbreak \iota_0 <\allowbreak \zeta_1$.\par\medskip
\noindent Initial window of $\sigma_{i+1}\sigma_i\sigma_{i+1}$. Columns: $(\alpha'_A,\zeta'_1),\allowbreak (\beta'_1,\zeta'_2),\allowbreak (\beta'_2,\alpha'_A),\allowbreak (\alpha'_B,\iota_0),\allowbreak (\theta_0,\beta'_1),\allowbreak (\beta'_3,\alpha'_B),\allowbreak (\alpha'_C,\beta'_3),\allowbreak (\zeta'_1,\beta'_2),\allowbreak (\zeta'_2,\alpha'_C)$.\\ Rows: $\theta_0 <\allowbreak \iota_0 <\allowbreak \alpha'_B <\allowbreak \zeta'_1 <\allowbreak \beta'_1 <\allowbreak \beta'_3 <\allowbreak \alpha'_A <\allowbreak \alpha'_C <\allowbreak \zeta'_2 <\allowbreak \beta'_2$.\par
\noindent Moves (18): $\mathsf{r}(\iota_0,\alpha'_B)$, $\mathsf{c}_{3}$, $\mathsf{d}^{X}(4,5;\alpha'_B\to\iota_0)$, $\mathsf{c}_{6}$, $\mathsf{d}^{O}(6,5;\alpha'_C\to\zeta'_2)$, $\mathsf{r}(\zeta'_1,\beta'_1)$, $\mathsf{r}(\iota_0,\beta'_1)$, $\mathsf{r}(\beta'_3,\alpha'_A)$, $\mathsf{d}^{O}(5,4;\beta'_3\to\zeta'_2)$, $\mathsf{c}_{2}$, $\mathsf{d}^{O}(2,1;\beta'_1\to\theta_0)$, $\mathsf{r}(\zeta'_1,\alpha'_A)$, $\mathsf{r}(\iota_0,\alpha'_A)$, $\mathsf{r}(\zeta'_1,\zeta'_2)$, $\mathsf{c}_{3}$, $\mathsf{d}^{O}(3,2;\beta'_2\to\zeta'_1)$, $\mathsf{c}_{2}$, $\mathsf{d}^{X}(2,1;\zeta'_2\to\iota_0)$.\par
\noindent Final window. Columns: $(\alpha'_A,\zeta'_1),\allowbreak (\theta_0,\iota_0),\allowbreak (\zeta'_1,\alpha'_A)$.\\ Rows: $\theta_0 <\allowbreak \alpha'_A <\allowbreak \iota_0 <\allowbreak \zeta'_1$.\par\medskip
\noindent Identification of the final windows: $\beta_0\leftrightarrow \alpha'_A,\allowbreak \\zeta_1\leftrightarrow \zeta'_1$.\par\bigskip
\subsection*{Case: degenerate ranks \texorpdfstring{$\{i,i+1,i+2\}$}{i,i+1,i+2}}
\noindent Initial window of $\sigma_i\sigma_{i+1}\sigma_i$. Columns: $(\alpha_A,\zeta_0),\allowbreak (\beta_0,\zeta_1),\allowbreak (\beta_1,\alpha_A),\allowbreak (\alpha_B,\beta_1),\allowbreak (\beta_2,\zeta_2),\allowbreak (\zeta_2,\alpha_B),\allowbreak (\alpha_C,\beta_0),\allowbreak (\zeta_0,\beta_2),\allowbreak (\zeta_1,\alpha_C)$.\\ Rows: $\zeta_0 <\allowbreak \beta_0 <\allowbreak \alpha_A <\allowbreak \alpha_C <\allowbreak \zeta_1 <\allowbreak \beta_1 <\allowbreak \beta_2 <\allowbreak \alpha_B <\allowbreak \zeta_2$.\par
\noindent Moves (8): $\mathsf{d}^{X}(5,4;\zeta_2\to\alpha_B)$, $\mathsf{d}^{O}(4,3;\alpha_B\to\beta_2)$, $\mathsf{d}^{O}(3,2;\beta_1\to\beta_2)$, $\mathsf{c}_{4}$, $\mathsf{d}^{O}(4,3;\alpha_C\to\zeta_1)$, $\mathsf{c}_{3}$, $\mathsf{r}(\zeta_1,\beta_2)$, $\mathsf{d}^{X}(2,3;\beta_2\to\alpha_A)$.\par
\noindent Final window. Columns: $(\alpha_A,\zeta_0),\allowbreak (\beta_0,\zeta_1),\allowbreak (\zeta_0,\alpha_A),\allowbreak (\zeta_1,\beta_0)$.\\ Rows: $\zeta_0 <\allowbreak \beta_0 <\allowbreak \alpha_A <\allowbreak \zeta_1$.\par\medskip
\noindent Initial window of $\sigma_{i+1}\sigma_i\sigma_{i+1}$. Columns: $(\alpha'_A,\zeta'_1),\allowbreak (\beta'_1,\zeta'_2),\allowbreak (\beta'_2,\alpha'_A),\allowbreak (\alpha'_B,\zeta'_0),\allowbreak (\zeta'_0,\beta'_1),\allowbreak (\beta'_3,\alpha'_B),\allowbreak (\alpha'_C,\beta'_3),\allowbreak (\zeta'_1,\beta'_2),\allowbreak (\zeta'_2,\alpha'_C)$.\\ Rows: $\zeta'_0 <\allowbreak \alpha'_B <\allowbreak \zeta'_1 <\allowbreak \beta'_1 <\allowbreak \beta'_3 <\allowbreak \alpha'_A <\allowbreak \alpha'_C <\allowbreak \zeta'_2 <\allowbreak \beta'_2$.\par
\noindent Moves (12): $\mathsf{d}^{O}(3,4;\zeta'_0\to\alpha'_B)$, $\mathsf{r}(\alpha'_A,\alpha'_C)$, $\mathsf{d}^{O}(5,4;\beta'_3\to\alpha'_C)$, $\mathsf{r}(\alpha'_C,\alpha'_A)$, $\mathsf{c}_{5}$, $\mathsf{d}^{O}(5,4;\alpha'_C\to\zeta'_2)$, $\mathsf{r}(\zeta'_1,\beta'_1)$, $\mathsf{d}^{X}(3,4;\alpha'_B\to\beta'_1)$, $\mathsf{r}(\beta'_1,\zeta'_1)$, $\mathsf{c}_{3}$, $\mathsf{r}(\zeta'_2,\beta'_2)$, $\mathsf{d}^{X}(2,3;\beta'_2\to\alpha'_A)$.\par
\noindent Final window. Columns: $(\alpha'_A,\zeta'_1),\allowbreak (\beta'_1,\zeta'_2),\allowbreak (\zeta'_1,\alpha'_A),\allowbreak (\zeta'_2,\beta'_1)$.\\ Rows: $\zeta'_1 <\allowbreak \beta'_1 <\allowbreak \alpha'_A <\allowbreak \zeta'_2$.\par\medskip
\noindent Identification of the final windows: $\zeta_0\leftrightarrow \zeta'_1,\allowbreak \\beta_0\leftrightarrow \beta'_1,\allowbreak \\alpha_A\leftrightarrow \alpha'_A,\allowbreak \\zeta_1\leftrightarrow \zeta'_2$.\par\bigskip

\end{document}
